在第一部分中，我们讨论语言建模与机器学习的基础。我们从语言如何被建模的基本原理出发，涉及词元与嵌入。接着讨论注意力这一概念——它既是 Transformer 架构的基石，也是大语言模型取得成功的根基。我们还会回顾两种重要的机器学习技术：监督学习与强化学习。在下一章中，我们将更具体地讨论一个典型训练流水线的各个步骤，从预训练、微调到推理。我们还会在该章讨论常见的基准测试。本部分最后一个内容章节围绕大语言模型中{\em 大}的那一面展开：我们讨论大语言模型的性能如何随参数数量、训练数据与算力而变化，并深入探讨那些既能受益于规模、又能把成本压到最低的高效模型。本部分末尾，我们将进入一个动手案例研究，你会在其中学到：从零打造你自己的小型大语言模型时，各个环节是如何串起来的。 
